Este Trabajo de Fin de Grado surge del interés por abordar uno de los retos más complejos y vigentes de la Ingeniería del Software contemporánea: el aseguramiento de la calidad y la evaluación sistemática de sistemas de inteligencia artificial generativa en contextos de alta sensibilidad, como es la educación.

El proyecto busca tender un puente riguroso entre los estándares internacionales de ingeniería del software (ISO/IEC 25010) y las necesidades didácticas reales de estudiantes y docentes, ofreciendo una metodología reproducible, fundamentada y contrastada experimentalmente.
